
\documentclass[10pt,a4paper]{article}
\usepackage[margin=17mm]{geometry}
\usepackage[T1]{fontenc}
\usepackage{lmodern}
\usepackage{amsmath,amssymb,booktabs,tabularx,array}
\usepackage{enumitem}
\setlength{\parindent}{0pt}
\setlength{\parskip}{4pt}
\setlist[itemize]{leftmargin=*,itemsep=3pt,topsep=3pt}
\newcommand{\block}[1]{\vspace{6pt}{\large\bfseries #1}\par\vspace{2pt}}
\newcommand{\E}{\mathbb{E}_s}
\pagestyle{empty}
\begin{document}
\begin{center}
{\LARGE\bfseries Model Formulations and Development}\\[5pt]
\end{center}

\block{1. Research Question}
\textbf{Which cities should the company enter, and how many vehicles should it allocate to each city?}
The underlying document develops three formulations: a basic model, a cost-minimization model, and a profit-maximization model.
City entry decisions $y_c$ and vehicle allocations $q_c$ are determined jointly before the demand scenario is realized. The same allocation applies across all scenarios.

\block{2. Core Structure}
\begin{tabularx}{\linewidth}{@{}>{\bfseries}p{22mm}X@{}}
\toprule
Decisions & $y_c\in\{0,1\}$: city entry; \quad $q_c\in\mathbb Z_{\ge0}$: number of vehicles \\
Inputs & $D_c^s$: scenario-specific demand; $\kappa_c$: service capacity per vehicle; $Q$: total fleet limit; $\bar q_c$: city-specific vehicle limit \\
Outcomes & $R_c^s=\min(D_c^s,\kappa_c q_c)$: served demand \\
 & $O_c^s=(\kappa_c q_c-D_c^s)^+$: unused service capacity \\
 & $U_c^s=(D_c^s-\kappa_c q_c)^+$: unmet demand \\
\bottomrule
\end{tabularx}
{\small Here, $(x)^+=\max(x,0)$. Once generated, demand scenarios are fixed inputs; demand is not chosen by the optimization model.}

\block{3. Comparison of the Three Objectives}
\begin{tabularx}{\linewidth}{@{}p{24mm}X@{}}
\toprule
\textbf{Model} & \textbf{Objective and interpretation} \\
\midrule
1. Basic & $\displaystyle \min\;\E\!\left[\sum_c(O_c^s+\rho U_c^s)\right]$ \\
 & Balances unused capacity and unmet demand. The parameter $\rho$ measures the relative importance of shortages. Entry and vehicle costs are excluded. \\
\addlinespace[5pt]
2. Cost & $\displaystyle \min\;\sum_c(F_cy_c+K_cq_c)+\E\!\left[\sum_c(c_c^+O_c^s+c_c^-U_c^s)\right]$ \\
 & Minimizes entry costs, vehicle costs, unused-capacity penalties, and unmet-demand losses. \textbf{This is the current baseline objective.} \\
\addlinespace[5pt]
3. Profit & $\displaystyle \max\;\E\!\left[\sum_c m_cR_c^s\right]-\sum_c(F_cy_c+K_cq_c)$ \\
 & Deducts entry and vehicle costs from the contribution margin on served trips. Here, $m_c$ is revenue per trip minus its variable cost. \\
\bottomrule
\end{tabularx}

\block{4. Common Constraints and Subsequent Revisions}
\[
y_c\le q_c\le\bar q_c y_c,\qquad
\sum_c q_c\le Q,\qquad
\sum_c y_c\le B,\qquad y_c\le a_c.
\]
Here, $a_c$ indicates whether operation in city $c$ is feasible.
The parameter $B$ limits the number of cities that can be entered simultaneously and is used only when such a restriction actually exists.

\block{5. Key Interpretation Notes}
\begin{itemize}
\item \textbf{Avoid double counting.} Do not include costs already covered by $K_c$ again in $c_c^+$. In the profit model, do not deduct the same forgone contribution margin again as an unmet-demand loss.
\item \textbf{Relationship between models.} Under identical demand and constraints, if $c_c^+=0$ and $c_c^-=m_c$, the cost and profit models yield the same optimal allocations. The cost objective equals a demand-dependent constant minus the profit objective.
\item \textbf{Unused capacity versus its penalty.} Unused capacity $O_c^s$ can be positive even when $c_c^+=0$. Only the additional unused-capacity penalty is zero; vehicle costs remain.
\item \textbf{Model development.} A fourth, multi-penalty formulation was introduced in the original outline. Following feedback on that formulation, Model 2 (cost minimization) was adopted as the baseline.
\end{itemize}
\end{document}
